\chapter{Allocation patterns in the renderer}
\label{ch:allocations}

\begin{aside}
The fastest allocation is the one that never happens. The
renderer's steady state is zero allocations per frame ---
because every byte of working storage was acquired once at
startup.
\end{aside}

\section{Where allocations live}

A frame's working storage:

\begin{itemize}[leftmargin=*]
  \item Front and back canvases: two \code{vector<uint64\_t>}.
  \item Emit buffer: one \code{string}.
  \item Per-element layout scratch: per-tree-walk
    \code{vector<LayoutInfo>}.
  \item Side-tables: \code{unordered\_map} for combining
    marks and link URLs.
\end{itemize}

All of these grow during the first few frames and stabilise.
\maya{} reserves typical sizes upfront so the growth happens
once.

\section{Reserving up front}

\begin{lstlisting}
front_.reserve(80 * 24);
back_.reserve(80 * 24);
emit_buf_.reserve(8192);
\end{lstlisting}

These are sized for typical TUI dimensions. A larger terminal
re-allocates; the cost is amortised.

\section{Pool allocation for elements}

The element tree is rebuilt each frame in a memory arena: a
linear bump allocator backed by a single large slab. The
slab is reset (pointer rewinds to start) between frames; no
free, no fragmentation.

\begin{lstlisting}
class Arena {
    char* begin_;
    char* cursor_;
    char* end_;
public:
    void* allocate(size_t n) {
        if (cursor_ + n > end_) grow();
        auto* p = cursor_;
        cursor_ += n;
        return p;
    }
    void reset() { cursor_ = begin_; }
};
\end{lstlisting}

A frame's element tree allocates from the arena; reset
between frames returns all memory at once.

\section{Move semantics}

Element construction returns by value. Without copies the
construction is cheap; \maya{}'s elements are typically
under 64 bytes and trivially copyable.

The DSL builders (\code{|}, \code{>>}) take rvalue references
and chain elements without copying. The result is one
element built in place.

\section{Avoiding short string allocations}

Small strings (under 16 bytes) live in the string's inline
storage thanks to SSO. \maya{}'s default \code{Style} names
and class names are short --- they don't allocate.

For longer strings, \code{string\_view} is used wherever
ownership is not needed.

\section{Per-frame churn}

What does allocate per frame:

\begin{itemize}[leftmargin=*]
  \item Element tree (arena, no malloc).
  \item Layout scratch (reused across frames).
  \item Emit buffer (reused, capacity stays).
\end{itemize}

What does not:

\begin{itemize}[leftmargin=*]
  \item Canvases (allocated once, resized only on terminal
    resize).
  \item Side-tables (cleared but storage retained).
  \item Signal graph (lives across frames).
\end{itemize}

\section{Measuring}

Tools like \code{heaptrack} or \code{jemalloc}'s profiling
show per-frame allocations. A well-tuned \maya{} app in
steady state shows zero \code{malloc} calls per frame after
warmup.

\begin{tryit}
Run your app under \code{heaptrack}. Type for 10 seconds.
Look at the allocations per second --- should be near zero
after the first few frames.
\end{tryit}

\section{The arena trade-off}

Bump allocation is fast but indiscriminate: anything in the
arena's lifetime persists until reset. Per-frame elements
fit perfectly; per-session state must not go in the arena.

\maya{} keeps signal storage, widget state, and config data
in regular allocators; only the element tree uses the arena.

\section{When to break the rules}

Per-element animation state with non-trivial destructors
(timers, callbacks) cannot live in the arena --- the reset
skips destructors. \maya{} uses regular allocation for these
and tolerates the cost; they are rare.

\section{Recap}

\begin{recap}
\begin{itemize}[leftmargin=*]
  \item Canvases, emit buffer reserved once.
  \item Element tree in a bump arena, reset per frame.
  \item Steady state allocates zero from the OS.
  \item String SSO handles short labels.
  \item Heavy state (signals, widget state) lives outside
    the arena.
\end{itemize}
\end{recap}

\section*{Workshop}

\begin{enumerate}[leftmargin=*]
  \item Profile your app with \code{heaptrack}. Identify
    the largest per-frame allocator.
  \item Move one of your widget's per-frame
    \code{vector<X>} into a reusable scratch buffer.
  \item Confirm by measurement that the per-frame
    allocation count drops.
\end{enumerate}
